\noindent \textit{\textcolor{red!60!black}{
First, I appreciate the improvement in measuring $d_{it}$ relative to all long-term obligations involving fixed contractual payments, but this only partly addresses my concerns. Preferred shares are not interest bearing, and thus may behave quite differently than debt during downturns or when firms are close to distress. I understand that there is nothing the author can do about this because bank debt is simply too small, but the paper should acknowledge this limitation more clearly and briefly discuss the potential confounders that cannot be addressed due to how firms structured their debt at that time.
}}

\bigskip

\noindent The referee raises a valid point, and we now address it explicitly in the manuscript. Preferred dividends are discretionary in a way that bond coupons are not---deferring them does not trigger default---and we agree that this distinction matters in distress, even though the flexibility was partial, as preferred issues of the era were typically cumulative. Because preferred shares enter the denominator of $\tilde{d}$ but not the numerator, the concern is that high-$\tilde{d}$ firms, which rely more on inflexible bond financing, may be more vulnerable to downturns generally, so that the 1933--1934 differential partly reflects debt rigidity rather than gold clause litigation.

\medskip
\noindent Two features of the evidence speak to this concern. The preferred-share placebo (column 6 of Table 4) replaces the numerator of $\tilde{d}$ with preferred shares and finds positive, insignificant 1933 and 1934 coefficients, whereas a general flexibility effect would tend to push these coefficients negative. And the temporal pattern (Figure 6) is concentrated in 1933--1934 and reverses after the 1935 ruling, which is difficult to reconcile with debt rigidity alone: that mechanism would also operate in the 1937--1938 recession, where we find no differential effect. The industry-year fixed effects and characteristic-interacted controls of Table 7 provide further protection against shocks correlated with industry and observable firm structure.

\medskip
\noindent These tests mitigate but cannot fully eliminate the concern: in this period virtually all corporate bonds carried gold clauses and preferred shares carried none, so gold exposure cannot be separated cleanly from liability composition. Section 5 now acknowledges this limitation explicitly.
